\setauthor{Elisa Kaltenberger}

\section{Gesellschaftliches Umfeld}
\setauthor{Elisa Kaltenberger}
Ein großer Teil der Bevölkerung nimmt alltägliche körperliche Aktivitäten als selbstverständlich wahr. Für Menschen mit körperlichen Einschränkungen ist diese Selbstverständlichkeit jedoch häufig nicht oder nur eingeschränkt gegeben. Barrieren, die für viele kaum wahrnehmbar sind, können für Betroffene nicht bewältigt werden. Dazu zählen unter anderem nicht abgesenkte Bordsteinkanten, unebene Pflastersteine oder zu niedrig platzierte Bankautomaten mit zu hohem Sichtschutz.

Die eingeschränkte Zugänglichkeit zu Infrastruktur und Dienstleistungen wirkt sich nicht ausschließlich auf  die Mobilität aus, sondern betrifft in hohem Maß die gleichberechtigte gesellschaftliche Teilhabe.  Selbstbestimmung, soziale Integration und berufliche Möglichkeiten hängen wesentlich davon ab, wie  barrierefrei öffentliche und private Räume gestaltet sind.

Mit steigender Lebenserwartung nimmt auch die Anzahl jener Personen zu, die temporär oder dauerhaft auf  Unterstützung angewiesen sind. Barrierefreiheit ist daher kein Randthema, sondern betrifft einen signifikanten Teil der Bevölkerung, siehe Abb. \ref{fig:lebensqualiteat}.

\section{Erfahrungsberichte aus Interviews}
\setauthor{Elisa Kaltenberger}
Die Berichte zeigen, dass Mobilitätseinschränkungen sowohl temporär als auch dauerhaft auftreten können. 
Besonders auffallend ist dabei nicht nur die physische Herausforderung, sondern auch die psychische. 
Gefühle wie Abhängigkeit, Unsicherheit oder Hilflosigkeit treten häufig auf, insbesondere wenn alltägliche 
Handlungen plötzlich nicht mehr selbstständig durchgeführt werden können.

Häufige Probleme die in den Interviews genannt wurden, sind:
\begin{itemize}
    \item Schwierigkeiten beim Überwinden von Stufen und Bordsteinkanten
    \item Probleme im öffentlichen Verkehr
    \item Unangenehmes fahren auf unebenem Untergrund (z.B. Pflastersteine)
    \item Abhängigkeit von fremder Hilfe
\end{itemize}

Gleichzeitig wurde deutlich, dass barrierefreie Maßnahmen wie Rampen, abgesenkte Übergänge oder gut 
zugängliche öffentliche Räume erheblich zur Selbstständigkeit und Lebensqualität beitragen können.

\section{Zielgruppe des Projekts}
\setauthor{Elisa Kaltenberger}

Die Zielgruppe des Projekts umfasst unterschiedliche Personengruppen, die direkt oder indirekt mit dem Thema 
Barrierefreiheit in Berührung kommen. 

Dazu zählen insbesondere Menschen, die in der Stadtplanung oder 
Architektur tätig sind, da sie maßgeblich an der Gestaltung öffentlicher Räume beteiligt sind und durch ein 
vertieftes Verständnis der Problematik barrierefreie Strukturen gezielter umsetzen können und somit mehr 
gleichberechtigte Teilhabe ermöglichen können.

Ebenso richtet sich die Simulation an Personen, die beruflich oder ehrenamtlich im Bereich der 
Barrierefreiheit arbeiten. Für diese bietet das Projekt eine zusätzliche Möglichkeit die Perspektiven 
anschaulich zu vermitteln und Sensibilisierungsarbeit praxisnah zu unterstützen.

Darüber hinaus gehören Angehörige von Menschen mit Mobilitätseinschränkungen zur Zielgruppe. Sie können 
durch die Simulation einen besseren Einblick in die alltäglichen Herausforderungen erhalten und dadurch 
mehr Verständnis für die Lebensrealität ihrer Familienmitglieder entwickeln und somit die Unterstützung 
gezielter gestalten und sich auf die Bedürfnisse der Betroffenen besser einstellen.

Eine weitere Zielgruppe sind alle, die nach einem Unfall kurz davor stehen, selbst auf einen Rollstuhl 
angewiesen zu sein. Für diese Personen bietet die Simulation die Möglichkeit, sich bereits vorab mit der 
neuen Lebenssituation auseinanderzusetzen und sich auf die Herausforderungen einzustellen. Dadurch kann 
die Simulation nicht nur zur Sensibilisierung beitragen, sondern auch als unterstützendes Werkzeug für die 
Vorbereitung auf die neue Lebensrealität dienen.

